This dissertation investigates how incentive structure shapes emergent strategy in multi-agent reinforcement learning, grounded in game theory. We ask a central question about what happens when we change the reward structure that agents optimise for. Specifically, we explore the journey from purely competitive to fully cooperative settings. How does this shift affect the strategies that agents learn in complex, stochastic environments?

We use a Formula 1 racing simulator to study this. Agents develop driving strategies through learning within this environment. Our experimental design isolates a single incentive variable, which we call alpha. This variable controls whether agents pursue their own success or their team's success. We sweep across the full spectrum from competitive to cooperative game structures. This allows us to measure how learned behaviour changes as incentives shift, holding everything else constant.

We test four state-of-the-art multi-agent reinforcement learning algorithms. We examine them across different levels of environmental randomness. Two critical findings emerge from this analysis. First, some algorithms suffer training instability when learning from experience accumulated over time. A problem emerges that is counterintuitive. More training data makes this instability worse, not better. 

Second, cooperation behaves in surprising ways as teams grow larger. With three agents working towards a shared goal, cooperation yields substantial gains. Joint performance improves by 83 per cent over purely competitive baselines. Yet when we scale to five agents, the result flips completely. Cooperation now degrades performance, regardless of how we structure incentives. We find a partial solution through curriculum scheduling, a training technique that separates two distinct failure modes. One failure mode is an ordering problem that curriculum scheduling can solve. The other is a credit assignment problem where agents struggle to connect their own actions to the outcomes they experience in large, interdependent teams.

The broader implication is striking. In multi-agent systems, incentive design matters as much as algorithm choice. As teams grow, standard architectural approaches simply break down. More data and better algorithms cannot overcome this alone. The future lies in learning architectures built specifically for cooperative multi-agent coordination at scale.
